% ============================================================
% PAGE 3 — Literature Review
% ============================================================

\section{Related Work}
\label{sec:related}

Object detection has a rich literature spanning several distinct
architectural paradigms. This section organizes prior work into
four thematic clusters most relevant to FusionDet.

\subsection{Region-Based Two-Stage Detectors}

The dominant paradigm through approximately 2018 was the
two-stage detector. The R-CNN family pioneered the use of selective
search to generate candidate bounding boxes, which were then
classified by a CNN. Faster R-CNN \cite{ren2015faster} replaced
selective search with a learned Region Proposal Network (RPN),
enabling end-to-end training and achieving significant speed
improvements over its predecessors. Mask R-CNN extended this framework to
instance segmentation by adding a parallel mask-prediction branch.
Despite high accuracy, two-stage detectors carry inherent latency
from the sequential proposal-then-classify pipeline, limiting their
applicability in real-time scenarios.

\subsection{One-Stage and Anchor-Free Detectors}

YOLO \cite{redmon2016yolo} fundamentally reframed detection as a
single regression problem: a single forward pass through a
fully-convolutional network directly predicts bounding boxes and
class probabilities on a spatial grid, achieving real-time inference.
Subsequent iterations---YOLOv3, YOLOv4 \cite{bochkovskiy2020yolov4},
and YOLOX \cite{ge2021yolox}---systematically improved the original
design through multi-scale prediction heads, better backbone choices,
and mosaic data augmentation. RetinaNet introduced Focal Loss
\cite{lin2017focal} to address the extreme foreground--background
class imbalance inherent in dense prediction, enabling a one-stage
model to match two-stage accuracy. Anchor-free approaches such as
FCOS and CenterPoint avoid the grid-search over anchor configurations
entirely, instead predicting object centers and regressing box
dimensions from those centers, simplifying training and reducing the
sensitivity to anchor design choices.

\subsection{Feature Pyramid and Multi-Scale Fusion}

Real-world scenes contain objects at vastly different scales, making
single-resolution prediction heads inadequate. Feature Pyramid
Networks \cite{lin2017feature} introduced a top-down pathway that
fuses high-resolution shallow features with semantically rich deep
features, becoming a near-universal component of modern detectors.
Path Aggregation Network \cite{liu2018path} augmented the FPN with a
bottom-up pathway, creating a bidirectional information flow that
benefits the deepest feature levels with recovered spatial detail.
BiFPN added learnable per-feature weights to the aggregation
junctions, enabling the network to selectively amplify the most
informative input at each merge point. NAS-FPN used neural
architecture search to discover an optimal pyramid topology,
demonstrating that hand-designed topologies may not be globally
optimal. FusionDet builds on the PANet topology but replaces
fixed-coefficient fusion with our trainable AFF module and
integrates channel attention at every pyramid level.

\subsection{Vision Transformers in Detection}

The self-attention mechanism \cite{vaswani2017attention} enables
transformers to model pairwise relationships between all tokens
simultaneously, providing a fundamentally global inductive bias
absent in convolutions. ViT \cite{dosovitskiy2021vit} demonstrated
that patch-based transformers trained on large datasets can
outperform CNNs on ImageNet classification. Swin Transformer
\cite{liu2021swin} introduced hierarchical feature maps and
shifted-window attention, recovering the multi-scale structure
essential for dense prediction tasks while bounding attention
complexity to be linear in the number of tokens. DETR
\cite{carion2020detr} replaced the detection head with a transformer
decoder that directly predicts a fixed set of object embeddings,
eliminating hand-designed anchor matching and NMS. Deformable DETR
accelerated convergence by restricting each query to attend to a
small set of key spatial locations, reducing the quadratic attention
cost to linear in practice.

A persistent challenge with transformer-based detectors is their
quadratic memory and compute scaling with spatial resolution, which
becomes prohibitive for high-resolution inputs. Linear-attention
reformulations \cite{katharopoulos2020transformers} address this by
factorizing the attention kernel, but typically sacrifice some
modeling power. FusionDet uses linear attention exclusively within
the transformer branch, preserving global context modeling at
manageable cost, and delegates fine-grained local feature extraction
to the parallel CNN branch where convolutions are provably more
efficient.

\subsection{Hybrid CNN--Transformer Architectures}

Several recent works have sought to combine CNN local feature
extraction with transformer global context. CoAtNet interleaves
depthwise convolution and self-attention stages, showing that
depthwise convolution subsumes relative-position attention in early
layers. CMT and MobileViT insert transformer blocks inside a
lightweight CNN backbone, achieving strong mobile-class performance.
DETReg and UP-DETR use unsupervised pretraining of the transformer
decoder to accelerate convergence on detection tasks. Unlike these
approaches---which treat the transformer as either an augmentation of
or a replacement for specific CNN stages---FusionDet maintains fully
independent parallel branches and unifies them through a dedicated
fusion module, preserving the full representational diversity of both
paradigms without architectural entanglement.
